\section{Motivation and Scope}
\label{sec_motivation}
The motivating scenario is a hybrid application whose managed heap is stable while
native allocations and resident pages grow.  A managed-only controller may see
little reason to collect, whereas a pressure-only controller reacts after the
system is already contended and may not identify which application's recent growth
is responsible.  Combining the signals could permit an earlier and more selective
request.

This paper evaluates that proposition only in simulation.  The four workload
generators model bursty camera, game-streaming, inference, and multitasking
allocation traces.  They do not model a named application, exact camera resolution,
TensorFlow Lite model, AR scene, vendor kernel, thermal state, or background-process
population.  No numerical observation in the evaluation should be interpreted as a
commercial-device measurement.

The scope restriction is important.  A deployable study must measure observer
error and delay, allocator survival and reuse, refaults during GC, writeback, energy,
thermal behavior, application kills, and long-run oscillation.  It must also
compare equal-reclaim-volume policies and publish raw traces and every per-run
record.  Section~\ref{sec_eval} lists these requirements alongside the simulation
results so that feasibility evidence is not confused with deployment evidence.
